\section{Research targets and conventions}
\label{poly:scope}

The newest incoming reports make the next exact questions unusually
specific. The higher-digamma report asks for a calculus of arbitrary
differential polynomials, for common shifted endpoints, and for explicit
sixth and seventh moments \cite{Higher}. The canonical fifth-moment report
asks which differential monomials survive integration by parts, with the
finite endpoint value retained, and conjectures that an optimal initial
set of harmonic-tail moments has full formal rank \cite{Fifth}.
The continuous-depth report asks for spectral derivatives at its odd
evaluation lattice \cite{Continuous}. These are the principal targets here.

The normalized Mittag--Leffler expansion for scalar powers, the
height-one Gamma integral, the centered Lerch reflection, and the original
two-tail moment formulas are antecedents. We reproduce the parts needed
to fix conventions and then prove their extensions. This distinction is
important because the repository already contains several independently
derived versions of related calculations \cite{Repo,Mixed,Tilted,Algebraic}.
The present audit is targeted, not an audit of every analytic assertion in
the collection.

\begin{center}
\small
\begin{tabular}{p{0.27\textwidth}p{0.62\textwidth}}
\toprule
Source question & Result and remaining scope\\
\midrule
Differential-polynomial moments & A complete residue and finite-part
theorem, with a finite reduction to convergent zeta coordinates for every
mixed monomial.\\[4pt]
Integration by parts at all weights & An explicit direct-sum normal form;
its dimension is a formal algebraic dimension, not independence of
the resulting numbers.\\[4pt]
Higher scalar Gamma normalization & Only the highest height-one Laurent
coordinate is singular; all remaining coordinates have convergent
representatives. Arithmetic reduction of those coordinates remains a
separate question.\\[4pt]
Odd-lattice Lerch spectral jets & Every order and mixed depth derivative,
with explicit local Gamma coefficients and convergent spectral evaluators.\\[4pt]
Optimal initial tail moments & Proved for all neighboring even orders
$p=q-2$. Arbitrary even gaps remain open.\\
\bottomrule
\end{tabular}
\end{center}

\subsection{Functions, finite parts, and spectral conventions}

Write $\psi=\Gamma'/\Gamma$ and $\psi^{(r)}=\dd^r\psi/\dd x^r$,
with $\psi^{(0)}=\psi$. Let
\[
 H_n^{(r)}=\sum_{k=1}^n k^{-r},\qquad H_n=H_n^{(1)},\qquad
 A_r(n)=\zeta(r)-H_n^{(r)}\quad(r\ge2).
\]
The empty sum is zero. A strict word has the convention
\[
 H_n(w_1,\ldots,w_d)=
 \sum_{n\ge k_1>\cdots>k_d\ge1}\prod_{j=1}^d k_j^{-w_j},
 \qquad H_n(\varnothing)=1.
\]
For convergent arguments, $Z(s_1,\ldots,s_d)$ is the corresponding
infinite strict multiple-zeta series. A derivative $\partial_j Z$
means differentiation in its $j$th spectral argument; its summand contains
$-\log k_j$. Classical definitions and normalizations of Gamma, Hurwitz
zeta, and Lerch's transcendent are those of the DLMF
\cite{DLMFGamma,DLMFZeta,DLMFLerch}.

The Stieltjes convention is
\begin{equation}\label{poly:stieltjes}
 \zeta(1+s,a)=\frac1s+
 \sum_{j\ge0}\frac{(-1)^j}{j!}\gamma_j(a)s^j,
 \qquad \re a>0.
\end{equation}
We write $\gamma_j=\gamma_j(1)$ and $\gamma=\gamma_0$.
The notation $[s]F(s)$ denotes a Laurent coefficient. It agrees with
$F'(0)$ only when $F$ is holomorphic at zero.

Suppose $f(x)=\sum_{j=1}^M p_jx^{-j}+O(1)$ near $0$. Our finite part is
\begin{equation}\label{poly:fp}
 \FP\int_0^1f(x)\dd x=
 \lim_{\varepsilon\downarrow0}
 \left\{\int_\varepsilon^1f(x)\dd x
 -\sum_{j=2}^M\frac{p_j\varepsilon^{1-j}}{j-1}
 +p_1\log\varepsilon\right\}.
\end{equation}
The coordinate is exactly $x$ and the upper endpoint exactly $1$.
Every finite-part statement below uses this definition. Normalized
primitives use principal logarithms on $\C\setminus(-\infty,0]$.

For a polynomial $P\in\C[u_0,u_1,\ldots,u_R]$, set
\begin{align*}
 f_P(x)&=P(\psi(x),\psi'(x),\ldots,\psi^{(R)}(x)),&
 \phi_P(Y)&=P(Y,0,\ldots,0),\\
 \cI(P)&=\FP\int_0^1f_P(x)\dd x.
\end{align*}
We call $P$ mixed when $\phi_P=0$. This is a precise algebraic condition:
each of its monomials contains at least one $u_r$ with $r\ge1$.
The word does not assert that the associated numerical periods are
independent. Give $u_r$ differential weight $r+1$.
